\documentclass[11pt]{article}
\ifdefined\pdfinfoomitdate\pdfinfoomitdate=1\fi
\ifdefined\pdftrailerid\pdftrailerid{}\fi
\usepackage[a4paper,hmargin=26mm,vmargin=25mm]{geometry}
\usepackage[T1]{fontenc}
\usepackage{lmodern,microtype,amsmath,amssymb,amsthm,mathtools}
\usepackage[hidelinks,bookmarksnumbered]{hyperref}
\hypersetup{pdftitle={Natural boundary and differential transcendence of the isotropic square-lattice Ising susceptibility},pdfauthor={Dehao Lin}}
\newtheorem*{theorem}{Theorem}
\newcommand{\ii}{\mathrm i}
\newcommand{\dd}{\,\mathrm d}
\newcommand{\C}{\mathbb C}
\newcommand{\R}{\mathbb R}
\newcommand{\Z}{\mathbb Z}
\newcommand{\T}{\mathbb T}
\DeclareMathOperator{\arsinh}{arsinh}
\setlength{\emergencystretch}{1.5em}
\raggedbottom
\title{Natural boundary and differential transcendence of the\\
isotropic square-lattice Ising susceptibility}
\author{Dehao Lin}\date{}
\makeatletter
\renewcommand{\@maketitle}{\begin{center}
  \vspace*{.5\baselineskip}{\LARGE\@title\par}\vspace{.5\baselineskip}
  {\large\@author\par}
  {\small School of Physics, Sun Yat-sen University\par
  \href{https://orcid.org/0009-0001-4551-8490}{ORCID 0009-0001-4551-8490}\par}
\end{center}}
\makeatother
\begin{document}
\maketitle
\begin{abstract}
For the zero-field isotropic Ising model on the square lattice, we prove
that the unit circle in $s=\sinh(2\beta J)$ is a natural boundary for the
full susceptibility in each phase. Neither susceptibility satisfies an algebraic
differential equation over $\C(s)$. At a dense set of boundary points
all singular contributions of the first resonant particle order have
the same phase.
The remaining finite-order collisions are integrable, and a Gaussian
bound in the lattice coordinates controls all sufficiently high orders.
The resulting divergence
also forces the separant of any proposed differential equation to vanish.
\end{abstract}

\section{The susceptibility in the two phases}

Guttmann and Enting found evidence of accumulating singularities in the
anisotropic high-temperature series \cite{GE}.
Nickel found singularities of the odd-particle integrals of the
isotropic square-lattice Ising susceptibility at a dense set of points
on the unit circle. He proposed that the circle is a natural boundary
of the full susceptibility \cite{Nickel}. The passage from a single
particle integral to the infinite sum requires control of cancellation.
Guttmann, Jensen, Maillard and Pantone asked whether the full
susceptibility satisfies an algebraic differential equation and tested
the series on both sides of the critical temperature \cite[Sec.~6]{GJMP}.

Let $J>0$ be the nearest-neighbor coupling, $\beta$ the inverse
temperature and $s=\sinh(2\beta J)$. With spins of modulus one, the
reduced zero-field susceptibilities are
\[
 X_H(s)=\sum_{a,b\in\Z}\langle\sigma_{00}\sigma_{ab}\rangle,\qquad
 X_L(s)=\sum_{a,b\in\Z}
 \bigl(\langle\sigma_{00}\sigma_{ab}\rangle_+-M^2\bigr),
 \qquad M^2=(1-s^{-4})^{1/4}.
\]
The first sum is taken at $0<s<1$, the second at $s>1$ in the plus
pure state. Spin reversal gives the same connected sum in the minus
state. Both sums include the origin. For a magnetic field in energy
units the physical susceptibilities are $\chi_\nu=\beta X_\nu$,
with the ordered-state response taken from positive field in the
plus phase \cite[(1.7)--(1.10)]{ONGP}.

\begin{theorem}
$X_H$ is holomorphic on $|s|<1$, with $X_H(0)=1$, and $X_L$ is
holomorphic on $|s|>1$. The unit circle is a natural boundary for
each. The same boundary holds for $\chi_H$ and for every local branch
of $\chi_L$ continued from $s>1$ in the exterior. At every point of
their respective real temperature intervals the germs of $X_\nu$ and
$\chi_\nu$ are differentially transcendental over $\C(s)$.
\end{theorem}

Here differential transcendence means that there is no nonzero
polynomial relation among the function and finitely many derivatives
with coefficients rational in $s$. We use the dispersion variable
$E=s+s^{-1}$. Both temperature domains map to the slit plane
$\C\setminus[-2,2]$, with $s=0$ corresponding to infinity. In the
series normalization of Guttmann and Maillard \cite[(4)--(5)]{GM},
the high- and low-temperature series $\widetilde\chi_H(w)$ and
$\widetilde\chi_L(w)$ are also differentially transcendental over
$\C(w)$, where $w=1/(2E)$.

\section{Site integrals and their sum}

For $q\in\T=\R/(2\pi\Z)$ let $z=e^{\ii q}$ and let $u=e^{-\gamma}$
be the root in the open unit disk of $u+u^{-1}=2(E-\cos q)$.
Write
\[
 K_{ij}=\frac{\sin((q_i-q_j)/2)}{\sinh((\gamma_i+\gamma_j)/2)},
 \qquad H(v,w)=\frac{v-w}{1-vw}.
\]
Subtracting the dispersion equations gives the single-valued square
\begin{equation}\label{pair}
 K_{ij}^2=-\frac{u_i u_j}{z_i z_j}
          \left(\frac{z_i-z_j}{1-u_i u_j}\right)^2
          = H(z_i,z_j)H(u_i,u_j).
\end{equation}
At $z_i z_j=1$ the middle expression gives the removable value of the product on the right.
For $N\ge1$, $A\in\Z_{\ge0}$ and $B\in\Z$, the real-temperature
spin-correlation expansion \cite[(3.13), (3.17)--(3.20), (4.1)--(4.2)]{ONGP}
gives the site coefficient
\begin{equation}\label{site}
 I_N(A,B;E)=\frac1{N!}\int_{\T^N}\prod_{i<j}K_{ij}^2
       \prod_i u_i^A z_i^B\,\frac{\dd q_i}{2\pi\sinh\gamma_i}.
\end{equation}
The ordered-state expansion also appears in \cite[Theorem~5.0]{PT}.
The connected correlation at $(a,b)$ is the sum of
$I_N(|a|,b;E)$ over positive odd $N$ in the high-temperature phase
and positive even $N$ in the low-temperature phase, multiplied by
$a_H=(1-s^4)^{1/4}/s$ and $a_L=(1-s^{-4})^{1/4}$, respectively.
The fourth roots tend to one at $s=0$ and $s=\infty$, respectively.

Write $D(v)=\prod_{i<j}H(v_i,v_j)^2$. The Cauchy determinant gives
$D(v)=\det[1+H(v_i,v_j)]$ since $1+H(v,w)=(1+v)(1-w)/(1-vw)$.
Clearing denominators in $\det[(1-v_iw_j)^{-1}]$ gives the two
Vandermonde factors by alternation and degree, with constant one by
induction at $v_Nw_N=1$.
We bound the pair product on compact angular contours where the roots
are holomorphic and $1-u_i u_j$ stays away from zero, including $i=j$.
At fixed $E$ replace every $q_i$ by $q_i+t$ and $u_i$ by $u(E,q_i+t)$.
Compactness gives a common $t$-disk. The complete product remains
holomorphic by \eqref{pair}, and its square is $D(z)D(u)$.
The root determinant has row norms at most $C\sqrt N$. In the angular
determinant each entry has bounded mean absolute value on the $t$-disk:
the possible poles of $(1-z_i z_j e^{2\ii t})^{-1}$ are simple,
with uniformly bounded residues, and $|t-t_0|^{-1}$ has bounded area
mean on a disk of fixed radius. Thus each angular row has mean norm
at most $CN$. Hadamard's inequality, the submean inequality for
log modulus and concavity of the logarithm therefore give
$|\prod_{i<j}K_{ij}^2|\le(C N^{3/4})^N$ in both parities.
After division by $N!$ this is summable for every $C$.

The decay in both lattice coordinates follows from their interchange.
For real $E>2$, on circles $|v_i|=|w_i|=\rho<1$ close to one, integrate
\[
 \frac1{N!}\prod_{i<j}H(v_i,v_j)H(w_i,w_j)
 \prod_i\frac{v_i^A w_i^B}
 {E-(v_i+v_i^{-1}+w_i+w_i^{-1})/2}
 \frac{\dd v_i\dd w_i}{(2\pi\ii)^2v_iw_i},\qquad A,B\in\Z_{\ge0}.
\]
This is the double-circle representation of \cite[Appendix~1]{TW}
in even parity. For either parity, each $v_i$ integral has only the
interior root residue $u_i^A/\sinh\gamma_i$. Formula \eqref{pair}
then allows the $w_i$ circles to expand to unit radius, giving
\eqref{site}. Exchanging the circles, reflecting all $q_i$, and
continuing in $E$ gives $I_N(A,B)=I_N(B,A)=I_N(B,-A)$ for $A,B\ge0$.
On each compact subset of the slit plane the real-angle $|u|$ is
bounded by some $\rho_0<1$ and the other factors are bounded. Hence
$|I_N(|a|,b)|\le (C N^{3/4})^N\rho_0^{N\max(|a|,|b|)}/N!$, summably
in both sites and orders.

For $\sigma=\pm1$ put $H_N^\sigma=\{q:e^{\ii\sum q_i}=\sigma\}$,
$Z=\prod_i u_i$ and $\dd\mu=\prod_{i<N}\dd q_i/(2\pi)$.
Summing \eqref{site} with weight $\sigma^{a+b}$ gives
\begin{equation}\label{bulk}
 C_N^\sigma(E)=\sum_{a,b\in\Z}\sigma^{a+b}I_N(|a|,b;E)
 =\frac1{N!}\int_{H_N^\sigma}\frac{1+\sigma Z}{1-\sigma Z}
          \prod_{i<j}K_{ij}^2\prod_i\frac1{\sinh\gamma_i}\dd\mu .
\end{equation}
The $a$ sum is geometric. The $b$ sum evaluates the absolutely
convergent Fourier series in the total angle at $e^{\ii\sum q_i}=\sigma$.
Thus $X_\nu=a_\nu F_\nu$, where $F_\nu(E)$ sums $C_N^+(E)$ in
the parity of the phase. Near $s=0$ the factors $u$, $1/\sinh\gamma$ and
$K_{ij}$ are uniformly $O(|s|)$, so $|C_N^+|\le3(C|s|)^{N^2}/N!$.
Since $C_1^+(E(s))=1/(E-2)=s/(1-s)^2$, the simple zero of $F_H(E(s))$
cancels the pole of $a_H$, giving $X_H(0)=1$.
Shifting all angles by $\pi$ and negating the roots gives
\begin{equation}\label{reflect}
 C_N^+(-E)=(-1)^N C_N^{(-1)^N}(E).
\end{equation}
In particular reflection in odd parity changes both signs in \eqref{bulk}.

\section{Contour bounds near the energy cut}

Fix $0<x<2$. At $E=x+\ii y$, $y>0$, write $\gamma=\ii\phi$ on
the physical lift. On the oscillatory interval $|r|<b=\arccos(x-1)$
the boundary value $g(r)=\arccos(x-\cos r)$ of $\phi$ satisfies
\[
 g'=-\frac{\sin r}{\sin g},\qquad
 g''=-\frac{x(1-\cos r\cos g)}{\sin^3g}<0 .
\]
Let $f$ be odd and periodic, equal to $r$ on $[-b,b]$ and linear
from $b$ to zero on $[b,\pi]$. For $c$ in a compact subinterval of
$(-b,b)$ put $\lambda=g'(c)$ and $q=r+\ii v(f(r)-c)$,
$-\pi\le r\le\pi$. Uniformly for small $v,y\ge0$,
\begin{equation}\label{phase}
 \Im(\lambda q-\phi(E,q))\ge a\{y+v(r-c)^2\},\qquad a>0 .
\end{equation}
In a compact interior interval the linear term at $y=0$ is
$v(r-c)(g'(c)-g'(r))\ge a_1v(r-c)^2$, with imaginary remainder
$O(v^3|r-c|^3)$. Increasing $y$ adds at least $a_2y$, since
$\Re(1/\sin\phi)>0$ there. Near $\pm b$, $\sin r(f(r)-c)>a_3>0$,
so $W=E-\cos q$ satisfies $\Im W\ge y+a_4v$. Here $G=-\Im\phi>0$
for $y+v>0$, and $\Im W/G=\sin(\Re\phi)\sinh(G)/G\to0$ at $W=1$.
Thus $G$ controls both the right side of \eqref{phase} and the
bounded multiple of $v$ contributed by $\lambda q$ there. On the
remaining exterior intervals $G>0$ already at $y=v=0$. These bounds
are locally uniform in $x$.

At $y=v=0$ the unit roots have arguments in $[-b,0]$. Two can be
reciprocal only at the branch value $1$. Near that value the preceding
deformation gives $|u|<1$ when $v+y>0$. By compactness the homotopy
therefore crosses no one-body or pair pole. At fixed small $v>0$,
their denominators have uniform gaps on the graph, also on a small
complex $E$-neighborhood.

\paragraph{The higher orders.}
At a nonzero site take $A=\max(|a|,|b|)$, $B=\min(|a|,|b|)$ and
choose $g'(c)=B/A$. Then $-\arccos(x/2)\le c\le0$. On the fixed graph
the exponential in \eqref{site} has modulus at most
$\exp[-a' A\sum_i(r_i-c)^2]$ by \eqref{phase}. Bound the pair product
by $(C N^{3/4})^N$ and integrate the Gaussian in each of the $N$ angles.
On an $E$-disk of radius $\delta/(1+A)$ the change in each $A\phi$
is bounded. Cauchy's estimate gives, including the origin by a fixed disk,
\begin{equation}\label{sites-bound}
 |\partial_E^j I_N(A,B;x+\ii y)|\le
 \frac{C_j^{\,N}N^{3N/4}}{N!}(1+A)^{j-N/2}.
\end{equation}
There are $8A$ sites on the shell of radius $A\ge1$. Consequently
the sum over $N>2j+4$ of either parity, with either sign in
\eqref{bulk}, has continuous derivatives through order $j$ up to
compact subintervals of $(0,2)$ from above.

\paragraph{The finite orders.}
The one-particle term is $C_1^\sigma=1/(E-2\sigma)$.
Fix $N\ge2$ and $\sigma$, and exclude the values $x=1+\cos\theta$ with
$e^{\ii N\theta}=\sigma$. On $H_N^\sigma$ take
$c=N^{-1}\sum_i f(r_i)$. The excluded values rule out the two constant
tuples at $\pm b$, so $c$ stays in a compact subinterval of $(-b,b)$.
On the graphs
$q_i=r_i+\ii v(f(r_i)-c)$ one has $\sum_i\Im q_i=0$, so the
angular constraint is preserved.
Along this homotopy \eqref{phase} gives
\[
 -\log |Z|\ge a\bigl(Ny+v\textstyle\sum_i(r_i-c)^2\bigr).
\]
For $y>0$ this keeps $1-\sigma Z$ nonzero throughout the homotopy.
Together with the one-body and pair gaps, Stokes' theorem on the
smooth pieces of the graph gives \eqref{bulk} on the final cycle.
The shared faces have opposite orientations.
At fixed small $v=\epsilon>0$, the only possible poles at $y=0$ are
the constant configurations satisfying
\begin{equation}\label{res}
 e^{\ii N q_0}=e^{\ii N g(q_0)}=\sigma,\qquad
 x=\cos q_0+\cos g(q_0),\qquad |q_0|<b .
\end{equation}
Away from neighborhoods of these tuples the fixed graph gives a
holomorphic integral across $x$. In each neighborhood write
$r_i=q_0+v_i$, $\sum_i v_i=0$. The graph is exactly
$q_i=q_0+L v_i$ with $L=1+\ii\epsilon$.

Put $\eta=E-x$, $D=1-\sigma Z$, $Q=\sum_i v_i^2$ and
$\Delta(v)=\prod_{i<j}(v_i-v_j)$. Write $\psi=g(q_0)$.
Since $\phi_E=-1/\sin\psi$ at the center and $\sum_i v_i=0$,
Taylor expansion gives
\begin{equation}\label{expansion}
 D=-\ii(d\eta+hL^2 Q)+O(|\eta|^2+|\eta|\,|v|+|v|^3),\qquad
 d=\frac{N}{\sin\psi}>0,\qquad h=-\frac{g''(q_0)}2>0 .
\end{equation}
For $\Im\eta\ge0$, $|d\eta+hL^2Q|\ge c(|\eta|+Q)$ because
$0<\arg L^2<\pi$. This estimate and the gap off the charts identify
the fixed-graph integral with the physical germ on an upper half-disk
by continuation from $\eta=\ii y$. Each squared pair supplies
$(v_i-v_j)^2$ times an analytic unit. The density in
$\prod_{i<N}\dd v_i$ is $\Delta^2 B(\eta,v)/D(\eta,v)$, where
$B(0,0)\ne0$ since the one-body factors, the graph Jacobian and
$1+\sigma Z$ are nonzero there. On the $(N-1)$-dimensional plane,
the radial power of $\Delta^2\dd v$ in $r=\sqrt Q$ is $N^2-2$.
Thus a pole $D^{-m}$ is uniformly integrable against this measure
whenever $N^2-2-2m>-1$.

\section{The first resonant order}

Take $\sigma=\pm1$ in odd parity and $\sigma=1$ in even parity.
Choose $0<\alpha<\pi/2$ with $e^{\ii N\alpha}=\sigma$ for some $N$
of the chosen parity. The mesh of these angular grids is $2\pi/N$,
so the energies $x=2\cos\alpha$ are dense in $(0,2)$ even after
omitting $x=1$. Since $g(\alpha)=\alpha$, the constant tuple at
$\alpha$ satisfies \eqref{res}. Let $n$ be the least resonant
particle number of this parity at $x$.

No such $x$ is an endpoint value $1+\cos\theta$ for a root-of-unity
angle $\theta$. Indeed $x=e^{\ii\alpha}+e^{-\ii\alpha}$ is an algebraic
integer. If $u=\cos\theta=x-1$ then $u$ is also an algebraic integer.
For $m\theta\in2\pi\Z$ its conjugates are roots of $T_m(X)-1$, where
$T_m(\cos t)=\cos(mt)$. All those roots lie in $[-1,1]$.
If any conjugate is $1$ or $-1$, irreducibility gives $u=1$ or
$u=-1$. Otherwise the product of the conjugates is an integer of
absolute value less than one, so $u=0$. These cases give $x=2,0,1$,
all excluded.

At each collision of order $n$, put $\psi=g(q_0)\in(0,\pi)$.
In the local density $\Delta^2 B/D$, the phase of $B(0,0)$ is the same
at all these collisions. Each divided
squared pair contributes $L^2/[4(\ii\sin\psi)^2]$, the one-body
factors contribute $(\ii\sin\psi)^{-n}$ and the Jacobian is
$L^{n-1}$. Since $1+\sigma Z=2$ at the center, this gives
\[
 B(0,0)\ \in\ \ii^{-n^2}L^{\,n^2-1}\R_{>0}.
\]

Set $k=\lceil(n^2-3)/2\rceil$ and take a common $\epsilon>0$ small
enough for all the finitely many orders $N\le2k+4$.
Approach $x$ along $\eta=L^2t$, $t\downarrow0$. At particle number $n$,
\eqref{expansion} gives $D=-\ii L^2(dt+hQ)+o(t+Q)$.
The radial power makes $D^{-j}$ uniformly integrable for $j\le k$.
In the $k$th derivative of $B/D$ the only
term with a higher pole is
\[
 (-1)^k k!\,B\,(\partial_\eta D)^k/D^{k+1}.
\]
Since $\partial_\eta D(0,0)=-\ii d$ and $d,h>0$, one constant rotation
works for all the finitely many charts of order $n$. In each sufficiently
small chart the rotated term has real part at least $c(t+Q)^{-k-1}$
for all small $t$. The angular integral
of $\Delta^2$ on $Q=1$ is positive. Restriction to
$\sqrt t<r=\sqrt Q<\delta$ therefore bounds its integral below by
a positive constant times
\[
 \int_{\sqrt t}^\delta r^{\,n^2-2-2(k+1)}\dd r .
\]
The exponent is $-1$ for odd $n$ and $-2$ for even $n$. In both
cases the contributions add to a divergent derivative, and all lower
derivatives stay bounded.
All lower orders continue across $x$. For $n<N\le2k+4$ there are
no endpoint tuples by the preceding endpoint test. At every resonance
of such an order even $D^{-k-1}$ is integrable, since
$N^2-2-2(k+1)>-1$. For $N>2k+4$ the site bound
\eqref{sites-bound} controls the entire remaining tail. Thus the
divergence holds for the full sum.

\section{Natural boundary and differential transcendence}

For positive $x$ use $\sigma=1$. For negative energies use
\eqref{reflect} with $\sigma=(-1)^n$. Complex conjugation supplies
the other side of the slit. The map $E=s+s^{-1}$ is locally
invertible at the resulting dense set on $|s|=1$. The factors
$a_\nu$ and every continued branch of $\arsinh(s)/(2J)$ are analytic
units there. By the chain and product rules the $k$th derivative of
each of $F_\nu$, $X_\nu$, and $\chi_\nu$ diverges along the corresponding
approach from within its domain, with all lower derivatives bounded. Density proves the
natural-boundary assertions, including at the omitted axes.

Suppose one of these functions, in the coordinate $z=E$ for $F_\nu$
or $z=s$ for $X_\nu$ or $\chi_\nu$, satisfies a differential
relation. After clearing denominators take $P(z,Y,\ldots,Y^{(r)})=0$
of least order and then least degree in its highest derivative.
The separant $S=P_{Y_r}(z,Y,\ldots,Y^{(r)})$ is not identically zero
by minimality. Continue the relation from the real interval to an
upper energy collar in $0<\Re E<2$, fixing the local inverse $s(E)$
and the physical branches. Delete
from $(0,2)$ the finitely many endpoint and two-cosine grid values
of orders $N\le2r+4$. On a remaining open subinterval
\eqref{sites-bound} and the fixed contours give continuous one-sided
derivatives through $r$, hence a continuous trace of $S$.

Each bounded range of particle numbers supplies only finitely many
two-cosine grid values in \eqref{res}. Thus the selected $x$ with
$k>r$ are dense in this subinterval. At each such point, differentiate the
relation $k-r$ times. This gives $S Y^{(k)}+R=0$, where $R$ is a
polynomial in $z$ and derivatives through order $k-1$. Along the
chosen ray $R$ stays bounded and $|Y^{(k)}|\to\infty$. Thus $S$
has trace zero, first on the dense set and then on the whole
subinterval. Schwarz reflection in $E$ across this zero trace and
the identity theorem force $S$ to vanish identically, a contradiction.
Finally, $X_\nu=a_\nu F_\nu$ gives
$\widetilde\chi_\nu(w)=F_\nu(1/(2w))$ in the normalization above.
A differential relation for either germ at $w=0$ would yield one
for $F_\nu(E)$ by the chain rule, since $w=1/(2E)$ is birational.

\begingroup\small
\begin{thebibliography}{9}
\setlength{\itemsep}{3pt}\setlength{\parsep}{0pt}
\bibitem{GE} A. J. Guttmann and I. G. Enting,
\emph{Solvability of some statistical mechanical systems},
Phys. Rev. Lett. \textbf{76}(3) (1996), 344--347.
\href{https://doi.org/10.1103/PhysRevLett.76.344}{doi:10.1103/PhysRevLett.76.344}.
\bibitem{Nickel} B. Nickel, \emph{On the singularity structure of the
2D Ising model susceptibility}, J. Phys. A \textbf{32}(21) (1999),
3889--3906. \href{https://doi.org/10.1088/0305-4470/32/21/303}{doi:10.1088/0305-4470/32/21/303}.
\bibitem{GJMP} A. J. Guttmann, I. Jensen, J.-M. Maillard and J. Pantone,
\emph{Is the full susceptibility of the square-lattice Ising model a
differentially algebraic function?}, J. Phys. A \textbf{49}(50) (2016),
504002. \href{https://doi.org/10.1088/1751-8113/49/50/504002}{doi:10.1088/1751-8113/49/50/504002}.
\bibitem{GM} A. J. Guttmann and J.-M. Maillard,
\emph{Automata and the susceptibility of the square lattice Ising model
modulo powers of primes}, J. Phys. A \textbf{48}(47) (2015), 474001.
\href{https://doi.org/10.1088/1751-8113/48/47/474001}{doi:10.1088/1751-8113/48/47/474001}.
\bibitem{ONGP} W. P. Orrick, B. Nickel, A. J. Guttmann and J. H. H. Perk,
\emph{The susceptibility of the square lattice Ising model: New developments},
J. Stat. Phys. \textbf{102}(3--4) (2001), 795--841.
\href{https://doi.org/10.1023/A:1004850919647}{doi:10.1023/A:1004850919647}.
\bibitem{PT} J. Palmer and C. Tracy, \emph{Two-dimensional Ising
correlations: Convergence of the scaling limit}, Adv. Appl. Math.
\textbf{2}(3) (1981), 329--388.
\href{https://doi.org/10.1016/0196-8858(81)90010-5}{doi:10.1016/0196-8858(81)90010-5}.
\bibitem{TW} C. A. Tracy and H. Widom, \emph{On the singularities in the
susceptibility expansion for the two-dimensional Ising model},
J. Stat. Phys. \textbf{156}(6) (2014), 1125--1135.
\href{https://doi.org/10.1007/s10955-014-1061-4}{doi:10.1007/s10955-014-1061-4}.
\end{thebibliography}
\endgroup
\end{document}
